\providecommand{\fig}[1]{\figurename~\ref{#1}}

\chapter{Introducción}
\glsresetall
\label{cap:Introduccion}
%gls{API}

%Este capítulo presenta el contexto general, la problemática, la justificación, el alcance y los objetivos del presente trabajo de titulación, el cual se orienta al diseño e implementación de un mecanismo de cifrado y autenticación transparente para proteger la comunicación entre el controlador de estación \gls{API} y el servidor \gls{SCADA} de la Micro-Red de la Universidad de Cuenca.
En este capítulo se presenta el contexto general del trabajo de titulación, orientado al diseño y análisis comparativo de métodos de magnificación de video para la detección de movimientos y vibraciones sutiles. Inicialmente, en la Sección~\ref{sec:Antecedentes} se desarrolla un recorrido crítico por la evolución doctrinal y tecnológica de las técnicas de magnificación de video; se repasan los enfoques pioneros y metodologías diseñadas para abordar mejoras en el procesamiento visual, analizando la transición hacia aproximaciones en fase y estrategias para la compensación de movimientos globales o reducción de ruido. A continuación, en la Sección~\ref{sec:IdentificacionProblema} se articula la identificación del problema de investigación a partir de las limitaciones del estado del arte. Posteriormente, en la Sección~\ref{sec:Justificacion} se expone la justificación del proyecto, resaltando su valor técnico en inspección no destructiva y la creación de un repositorio de datos libre para la comunidad académica. Luego, en la Sección~\ref{sec:Alcance} se delimitan las fronteras del trabajo, especificando los escenarios de prueba sintéticos y reales, los métodos evaluados y las condiciones de experimentación. Finalmente, en la Sección~\ref{sec:Objetivos} se formalizan el objetivo general y los objetivos específicos que dirigen la ejecución metodológica del estudio.

\section{Antecedentes}
\label{sec:Antecedentes}
La detección y caracterización de movimientos, así como de vibraciones sutiles, constituye un desafío en múltiples aplicaciones científicas e industriales debido a que gran parte de la información asociada a estos fenómenos presenta amplitudes inferiores a la capacidad de percepción del sistema visual humano e incluso a la resolución de análisis proporcionada por métodos convencionales de inspección visual. Esta limitación dificulta la identificación de variaciones periódicas, micro-desplazamientos y cambios dinámicos presentes en estructuras, componentes mecánicos y sistemas biológicos que contienen información relevante para la evaluación de su estado y comportamiento.\\


Ante este desafío, la evolución doctrinal y tecnológica de la visión por computadora ha impulsado el desarrollo de técnicas computacionales capaces de aislar e intensificar tales variaciones en secuencias de video. Inicialmente, el análisis de micro-movimientos se abordó desde la \textbf{perspectiva Lagrangiana} 

\section{Identificación del Problema}
\label{sec:IdentificacionProblema}

La detección y caracterización de movimientos y vibraciones sutiles constituye un desafío en múltiples aplicaciones científicas e industriales debido a que gran parte de la información asociada a estos fenómenos presenta amplitudes inferiores a la capacidad de percepción del sistema visual humano e incluso a la resolución de análisis proporcionada por métodos convencionales de inspección visual. Esta limitación dificulta la identificación de variaciones periódicas, microdesplazamientos y cambios dinámicos presentes en estructuras, componentes mecánicos y sistemas biológicos, los cuales contienen información relevante para la evaluación de su estado y comportamiento.\\

Las técnicas de magnificación de video han permitido abordar este problema mediante la amplificación computacional de variaciones temporales presentes en secuencias de video, dando origen a diferentes enfoques como la magnificación Euleriana \cite{Wu2012}, la magnificación basada en fase \cite{Wadhwa2013} y la magnificación mediante pirámides de Riesz \cite{Wadhwa2014}. Cada uno de estos métodos presenta ventajas y limitaciones relacionadas con la calidad de la magnificación, la sensibilidad al ruido, la aparición de artefactos visuales, la precisión en la representación de movimientos sub-píxel y el costo computacional requerido para su implementación. Como consecuencia, en los últimos años se han desarrollado estrategias complementarias orientadas a mejorar el desempeño de estos algoritmos mediante técnicas robustas frente a grandes movimientos, filtrado adaptativo y procesamiento avanzado de señales \cite{Elgharib2015,Takeda2018,Takeda2022}.\\

No obstante, los trabajos reportados en la literatura se han orientado principalmente al desarrollo de nuevos algoritmos o a la optimización de métodos existentes para aplicaciones específicas, tales como la estimación de parámetros fisiológicos, el reconocimiento de microexpresiones, el monitoreo estructural o la detección de movimientos imperceptibles. Asimismo, los estudios comparativos disponibles suelen restringirse al análisis de un número reducido de indicadores de desempeño, como el tiempo de ejecución, el factor de amplificación o métricas de calidad de imagen, sin establecer un marco experimental unificado que permita evaluar simultáneamente aspectos relacionados con la fidelidad de la magnificación, la estabilidad temporal, la precisión en la estimación de frecuencias, la aparición de artefactos visuales y el rendimiento computacional de los diferentes enfoques \cite{Nandana2020}. Adicionalmente, existe una limitada caracterización del impacto que tiene la frecuencia de adquisición de las secuencias de video sobre el desempeño de estos algoritmos, particularmente cuando se dispone de cámaras de alta velocidad capaces de registrar fenómenos dinámicos de mayor frecuencia. Esta situación dificulta establecer criterios objetivos para determinar el método de magnificación más apropiado de acuerdo con las condiciones de captura, los requerimientos computacionales y la naturaleza del fenómeno analizado.\\

En consecuencia, se identifica la necesidad de desarrollar un sistema de procesamiento de video que permita implementar y evaluar, bajo un mismo marco experimental, los principales enfoques de magnificación de video utilizando secuencias sintéticas, bases de datos de referencia y secuencias adquiridas mediante una cámara de alta velocidad. Este análisis permitirá caracterizar de manera objetiva sus ventajas, limitaciones y condiciones de operación mediante métricas cuantitativas y cualitativas, contribuyendo a establecer criterios técnicos para la selección del método más adecuado según la aplicación y favoreciendo, de manera complementaria, la generación de un repositorio de datos experimentales que pueda ser utilizado en futuras investigaciones relacionadas con la magnificación de video y el análisis de movimientos sutiles.


\section{Justificación}
\label{sec:Justificacion}



\section{Alcance}
\label{sec:Alcance}
El presente trabajo se restringe al análisis comparativo y caracterización de tres enfoques representativos de magnificación de video: el método Euleriano basado en filtrado temporal y pirámides Laplacianas \cite{Wu2012}, el método basado en fase mediante pirámides orientables \cite{Wadhwa2013} y el método basado en pirámides de Riesz \cite{Wadhwa2014}. La selección de estos enfoques responde a que constituyen la base de la mayoría de las técnicas desarrolladas posteriormente en la literatura y representan las principales estrategias empleadas para la amplificación de movimientos y vibraciones sutiles.\\

Para ello, se desarrollará una arquitectura modular en \textit{Matlab} que integre las etapas de adquisición, preprocesamiento, descomposición, filtrado, amplificación y reconstrucción, permitiendo implementar cada algoritmo bajo un mismo entorno de procesamiento. Esta estructura facilitará la evaluación individual de cada método y garantizará que las comparaciones se realicen bajo condiciones homogéneas, evitando que diferencias en la implementación influyan sobre los resultados obtenidos.\\

La validación experimental se realizará mediante secuencias sintéticas y secuencias reales provenientes de bases de datos de referencia, entre ellas el conjunto de videos del MIT CSAIL \cite{Wu2012} y el dataset CASME II \cite{Yan2014}, las cuales serán complementadas con secuencias capturadas mediante una cámara de alta velocidad de hasta 1000 cuadros por segundo en escenarios controlados. A partir de estas pruebas se analizará el desempeño de los algoritmos considerando métricas relacionadas con la calidad de la magnificación, la estimación de frecuencias dominantes, la estabilidad temporal, la aparición de artefactos visuales y el rendimiento computacional, evaluando además su factibilidad de ejecución en modalidades \textit{offline}, cuasi-tiempo real y tiempo real.\\

Como resultado, se obtendrá una caracterización comparativa de los algoritmos evaluados, identificando sus ventajas, limitaciones y condiciones de operación, con el propósito de establecer criterios técnicos para la selección del método más adecuado según la aplicación y los requerimientos de procesamiento. De manera complementaria, las secuencias adquiridas durante el desarrollo experimental permitirán consolidar un conjunto de datos orientado a la evaluación de algoritmos de magnificación de video, el cual podrá ser divulgado con fines académicos para favorecer la reproducibilidad de futuras investigaciones.\\

El alcance de esta investigación no contempla el desarrollo de nuevos algoritmos de magnificación de video ni la optimización de los métodos seleccionados mediante técnicas de aprendizaje automático u otras estrategias de mejora. Asimismo, la evaluación se limitará a escenarios experimentales controlados y a los conjuntos de datos considerados en el estudio, por lo que la validación en entornos operacionales específicos quedará como una línea de trabajo futura.





\section{Objetivos}
\label{sec:Objetivos}

\subsection{Objetivo general}
Diseñar y comparar métodos de magnificación de video basados en enfoques Eulerianos, de fase y técnicas avanzadas mediante su implementación en un sistema de procesamiento de video con el fin de identificar sus ventajas y limitaciones de su operación en la detección de movimientos y vibraciones sutiles en entornos controlados.

\subsection{Objetivos específicos}
\begin{itemize}
    \item Analizar métodos de magnificación de video basados en enfoques Eulerianos, basados en fase y pirámides de Riesz, identificando sus principios de funcionamiento, parámetros de configuración y condiciones de aplicación para sustentar su evaluación comparativa.
		\item Desarrollar en MATLAB la arquitectura funcional de un sistema de procesamiento de video que integre las etapas de adquisición, preprocesamiento, descomposición, amplificación y reconstrucción para el análisis de variaciones sutiles en secuencias de video.
		\item Establecer escenarios de prueba en entornos controlados, integrando el uso de una cámara de alta velocidad, orientados a la captura y análisis de secuencias de video con variaciones sutiles bajo condiciones experimentales representativas para la validación empírica del sistema de procesamiento y la consolidación de un repositorio de datos de libre acceso para el ámbito académico.
		\item Caracterizar el desempeño de los algoritmos de magnificación de video mediante secuencias sintéticas y reales, empleado métricas cuantitativas y cualitativas para la identificación de ventajas, limitaciones y condiciones de operación considerando su rendimiento computacional, así como su factibilidad en escenarios offline, cuasi-real y tiempo real.
\end{itemize}